% claim-ref: kuhli-loach#breeding-evidence
% claim-ref: kuhli-loach#sex-study
% claim-ref: kuhli-loach#observation-sequence
% claim-ref: kuhli-loach#community-choice
% claim-ref: kuhli-loach#shrimp-risk
% claim-ref: kuhli-loach#nursery-separation
% claim-ref: kuhli-loach#hatch-interval
% claim-ref: kuhli-loach#maturity-age
\CompanionHeader{REPRODUCTION}{Kuhli loaches}{Pangio spp. / observation and compatibility}{A reproductive-care reference, not plant propagation or a validated home spawning recipe.}
\Context{Breeding is uncommon in trade-group advice. P. semicincta reports remain poorly documented in the specialist profile. No reliable hatch interval, maturity age or induction procedure is established for these fish. \textbf{[KUH-COOP, KUH-SF]}}
\Step{1}{Establish what is known}{Keep identity provisional and document ordinary health and feeding. A full abdomen alone does not identify sex or prove eggs. Budi et al.'s study of 50 fish labeled P. kuhlii found body and pectoral-fin differences; its abstract is not a diagnostic key for this pair. \textbf{[KUH-SEX]}}
\Step{2}{Plan before a breeding attempt}{Proposed: first arrange species-appropriate nursery conditions, suitable first food, filtration and a rehoming plan. No routine home protocol is validated here. Avoid deliberate overfeeding, hormone use, stripping or sudden chemistry/temperature changes. \textbf{[KUH-VET, KUH-SEX; proposed]}}
\Step{3}{Record distinct observations}{Proposed: if eggs or young appear, record dates, photographs, water and actual feeding. Distinguish eggs seen, hatching, free movement, first feeding and continuing growth. Obtain species-specific help before intervention; do not dismantle the occupied tank to search. \textbf{[KUH-SEX, KUH-VET; proposed]}}
\Panel{Tankmates / candidates and exclusions}{Small peaceful rasboras or tetras are candidates only after species-specific water and food-competition review. Aggressive, nipping or predatory fish, and incompatible-temperature species, are excluded from this proposed plan. No fish is guaranteed compatible merely by a family name. \textbf{[KUH-COOP, KUH-SF; proposed]}}
\Panel{Cherry shrimp / adult coexistence is not juvenile safety}{Keeper reports describe peaceful Kuhlis, but their small-prey diet and baby shrimp vulnerability justify caution. Adults may coexist; tiny juveniles, eggs and freshly molted shrimp can remain vulnerable. Cover cannot guarantee recruitment. No measured Kuhli predation rate was found. \textbf{[KUH-SHRIMP, KUH-SF, KUH-COOP; inference]}}
\Panel{Keep the future nursery plan separate}{Daniel's possible 10-gallon cherry nursery and later adult transfers are for a separate shrimp plan. It is not automatically a Kuhli breeding or quarantine system. Review capacity, chemistry and separate equipment before assigning any of these roles. \textbf{[KUH-VET, KUH-SHRIMP; proposed]}}
\Panel{Observation record}{Date / event: \hrulefill\par
Water / feeding / next check: \hrulefill}
\Sources{Sources (clickable): \href{https://scholar.unair.ac.id/en/publications/primary-and-secondary-sexual-characteristics-of-kuhli-loach-pangi/}{KUH-SEX: Budi et al. 2024 abstract}; \href{https://www.aquariumcoop.com/blogs/aquarium/care-guide-for-kuhli-loaches}{KUH-COOP: Co-Op Kuhli guide}; \href{https://www.seriouslyfish.com/species/pangio-semicincta/}{KUH-SF: Seriously Fish species profile}; \href{https://www.msdvetmanual.com/exotic-and-laboratory-animals/aquarium-fish/management-of-aquarium-fish}{KUH-VET: MSD veterinary care}; \href{https://www.aquariumcoop.com/blogs/aquarium/cherry-shrimp-care}{KUH-SHRIMP: Co-Op cherry shrimp}. Full titles, scopes and claims: adjacent YAML worksheets. Published care, proposed plans and unknowns remain distinct.}{kuhli-loach / reproduction}
